\section{Retrieval Harness}
\label{sec:harness}

\paragraph{Device and runtime.}
All measurements are taken on an Apple~M2 (Mac14,2, 8-core GPU) with MLX
\citep{mlx}; the retrieval models run in fp16, while the ingest model of
Section~\ref{sec:distill} runs in bfloat16. The vector index is a local Qdrant
instance without a server process \citep{qdrant}; it stores 3{,}134 points with
fp32 vectors.

\paragraph{Models.}
Dense retrieval uses \texttt{Qwen3-Embedding-0.6B} with 1{,}024-dimensional
output and a maximum input length of 2{,}048 tokens, with the model's query
instruction prepended to each query \citep{qwen3embedding}. Reranking uses
\texttt{Qwen3-Reranker-0.6B} from the same family. Candidates are length-bucketed
and scored in batches of up to sixteen pairs, and the reranker sustains about
1.2 pairs per second end to end at the actual chunk lengths.

\paragraph{Sparse channel and fusion.}
The default sparse channel is an in-memory BM25Okapi implementation with a
Russian Snowball stemmer \citep{robertson2009bm25,porter1980algorithm} and
corpus-calibrated average length. We additionally evaluate the
\texttt{Qdrant/bm25} sparse backend, which applies IDF on the server side; the
backend is eighteen stop-word files and a Python implementation, with no neural
network involved. Dense and sparse lists are combined with reciprocal rank fusion
\citep{cormack2009reciprocal} with the rank constant fixed at $k=60$.

\paragraph{Metrics.}
We report recall at 1, 5, 10 and 50, MRR@10, nDCG@10, document-level recall@10
and the fraction of queries for which the gold group is absent from the top 50
(\emph{miss@50}); recall@$k$ appears as R@$k$ in the tables. All metrics use the
gold groups of Section~\ref{sec:queries}. Differences between arms are reported
with paired bootstrap intervals over queries (20{,}000 resamples, fixed seed).
Latency is measured per query after warm-up and reported as the median and the
95th percentile over the evaluated query set; arms \arm{c} and \arm{d} are timed
on the sixty-query subset of Section~\ref{sec:results-main}.

\paragraph{Arms.}
Table~\ref{tab:arms} lists the eight configurations. Four of them form a
$2\times2$ design over the index content (flat chunk or contextualized chunk)
and the presence of the reranker; the remaining four isolate dense-only
retrieval and the distilled context. The two indexes differ only in the text
that is embedded, so a difference between arm~\arm{a} and arm~\arm{b} is
attributable to the context alone.

\begin{table}[t]
\centering
\small
\caption{Retrieval arms. \emph{Index text} is what the embedder sees;
\emph{query path} is what happens after the first stage. Arms \arm{c} and
\arm{d} rerank a pool of sixty-four candidates with a 0.6B cross-encoder.}
\label{tab:arms}
\begin{tabular}{lll}
\toprule
Arm & Index text & Query path \\
\midrule
\arm{a} & flat chunk & dense, BM25, RRF \\
\arm{b} & context $\Vert$ chunk & dense, BM25, RRF \\
\arm{c} & context $\Vert$ chunk & dense, BM25, RRF, reranked \\
\arm{d} & flat chunk & dense, BM25, RRF, reranked \\
\arm{e} & flat chunk & dense only \\
\arm{f} & context $\Vert$ chunk & dense only \\
\arm{g} & distilled context $\Vert$ chunk & dense only \\
\arm{h} & distilled context $\Vert$ chunk & dense, BM25, RRF \\
\bottomrule
\end{tabular}
\end{table}
